\section{Benchmark}\label{sec:bench}
\paragraph{Models.}
\begin{itemize}[nosep]
\item \textbf{ECFP models.} Count-based Morgan fingerprints (radius 2, 2{,}048 bits, RDKit) with histogram gradient boosting (GBM), random forest (RF, 300 trees) or a multilayer perceptron (MLP), all from scikit-learn. ECFP+GBM is the reference model for paired comparisons.
\item \textbf{Frozen embeddings.} ChemBERTa-77M-MLM~\cite{ahmad2022chemberta2} (DeepChem/ChemBERTa-77M-MLM) and MoLFormer-XL trained on 10\% of its corpus~\cite{ross2022molformer} (ibm-research/MoLFormer-XL-both-10pct). Embeddings are the attention-masked mean of the last hidden layer for the canonical SMILES (maximum 256 tokens), with a GBM head. Hugging Face revisions are pinned: ChemBERTa \texttt{ed8a537}, MoLFormer \texttt{7b12d94} (March 2024; the latest revision requires a newer \texttt{transformers} release).
\item \textbf{D-MPNN.} A directed message-passing network~\cite{yang2019analyzing} implemented in PyTorch (mean pooling, early stopping on 10\% of the training set, at most 60 epochs, patience 10). Run on five tasks with three seeds because of its cost on CPU.
\item \textbf{Fine-tuned ChemBERTa.} ChemBERTa-77M-MLM fine-tuned end to end with a linear head on the mean-pooled representation (AdamW, learning rate for the encoder chosen from $\{2\cdot10^{-5},10^{-4}\}$ on an inner validation split, $10^{-3}$ for the head, batch size 32, early stopping with patience 5, at most 40 epochs, standardised targets). Run on the same five tasks with three seeds (five on one task).
\end{itemize}

\paragraph{Hyperparameters.}
All ECFP models, embedding heads and the D-MPNN receive the same budget: three configurations sampled from a fixed grid (Appendix~\ref{app:repro}), selected by RMSE on a random 20\% of the training set, then refitted on the full training set. The inner validation split is random for all outer splits, i.e.\ it does not mimic the scaffold or temporal shift. This is a small budget; it is equal across these families. The fine-tuned ChemBERTa has a different budget (two learning rates, 10\% validation, no refit). ECFP+MLP (standardised count fingerprints, scikit-learn MLP) performs poorly, below the training mean on temporal splits; it should be read as an untuned baseline, not as representative of neural networks on fingerprints.

\paragraph{Splits and seeds.}
Test sets are 20\% of the molecules. Random split: uniform at random. Scaffold split: Bemis--Murcko scaffold groups in random order are assigned to the test set while they fit, so large scaffold groups stay in training. Temporal split~\cite{sheridan2013timesplit}: the 20\% of molecules first reported most recently form the test set (deterministic; the test sets start between 2010 and 2022 depending on task, median 2019). Molecules first reported in the cut-off year are divided between training and test by file order, and training labels average all documents of a molecule, including documents published after the cut-off; the split is therefore slightly optimistic. Random and scaffold splits use five seeds; for the temporal split the seeds affect only model training.

\paragraph{Evaluation.}
For each task, split, seed and model we compute RMSE, $R^2$ and $f$ on the test set, average over seeds, and then over tasks. Uncertainty for averages is a bootstrap over test molecules (300 resamples). Comparisons with ECFP+GBM are paired by task: we report the mean difference in $f$, the number of tasks on which the model is better, and a Wilcoxon signed-rank test over the 17 tasks with Holm correction over all model--split comparisons. The test treats tasks as independent, which is approximate: CHEMBL237 appears with both $K_i$ and $\mathrm{IC}_{50}$, and related receptors (D$_2$/D$_3$, $\mu$/$\kappa$ opioid) share ligands. Models run on a subset of tasks are compared only on that subset and without tests.
